% !TEX root =../main.tex

\subsection*{Part A: PVSTP}

Con la red montada, se conecto un computador ajeno al laboratorio y se empezo a
configurar uno a uno los switches, los comandos utilizados en general se presentan a 
continuacion.

\begin{darkcode}[Switch-X: Configuracion basica]
Switch-Y>ena
Switch-Y#config t
'
Enter configuration commands, one per line.  End with CNTL/Z.
'
Switch-Y(config)#hostname Switch-X
Switch-X(config)#spanning-tree mode pvst
Switch-X(config)# exit
Switch-X#show c
Switch-X#show spanning-tree 
'
VLAN0001
  Spanning tree enabled protocol ieee
  Root ID    Priority    32769
             Address     \swdA
             Cost        19
             Port        3 (FastEthernet0/3)
             Hello Time   2 sec  Max Age 20 sec  Forward Delay 15 sec

  Bridge ID  Priority    32769  (priority 32768 sys-id-ext 1)
             Address     0026.0af1.6b00
             Hello Time   2 sec  Max Age 20 sec  Forward Delay 15 sec
             Aging Time  15  sec

Interface           Role Sts Cost      Prio.Nbr Type
------------------- ---- --- --------- -------- ------
Fa0/1               Desg FWD 19        128.1    P2p 
Fa0/2               Altn BLK 19        128.2    P2p 
Fa0/3               Root FWD 19        128.3    P2p 
Gi0/1               Altn BLK 19        128.25   P2p 
'
\end{darkcode}

Siendo este el ultimo mensaje del que se determino la informacion mas relevante. Pero
tambien donde se encontraron los primeros problemas. En el ejemplo anterior se hace la
configuracion basica del dispositivo: entrar a modo privilegiado y cambiar el nombre
del switch. Luego se corre el comando \texttt{spanning-tree mode pvst} para entrar en
modo PVSTP.

Entonces se salio del modo de configuracion global para comprobar que se configuro
y conecto el switch usando el comand \texttt{show spanning-tree}. En el ejemplo se
muestran los resultados the \texttt{show spanning-tree} del switch B. Estos dicen que
en el puerto \texttt{Fa0/1} el designado---el computador---, el
\texttt{Fa0/2} un alternativo,  en \texttt{Gi0/1} otro alternativo y en \texttt{Fa0/3}
el root--el switch D---.

Si se obervan los mismo valores en otros switches se ve que todos apuntan al switch D
como el root y otros switches que no son el D como alternativos. Mientras que la PC
simpre es designado.

Pero ese mensaje no fue el originalmente recibido al primer intento, cuando se
hicieronn las conexiones se encontraron multiples cables que no funcionaban
adecuadamente de los que no pasaban las señales, la conexion \texttt{Fa0/1} no la podia
\textit{ver} el switch. Por tanto, por cada uno de los switch que le faltaba
informacion, se reviso el cable que impedia la conexion usando un prober y se cambio si
el prober indicaba un daño fisico.

Cabe anotar el hecho de que si la conexion al switch root D estaba dañada, el
resultado era tomar un switch alternativo y cambiarlo a root.

Una vez se arreglaron todas las conexiones, se corrio el mismo ultimo comando que en
el ejemplo anterior,
\texttt{show spanning-tree}, y se dejo en la tabla \ref{tab:spanning-tree-before} la
informacion recolectada y se dejo aparte en la table
\ref{tab:spanning-tree-roles-before} los roles de cada switch acorde a los otros
switch.

\begin{center}
  \begin{tabular}{ l c c c c }
    \hline
    --- & Switch-A & Switch-B & Switch-C & Switch-D \\
    \hline
    MAC & \swaA & \swbA & \swcA & \swdA \\
    BID MAC & \swaA & \swbA & \swcA & \swdA \\
    BID Priority & 32769 & 32769 & 32769 & 32769 \\
    Root bridge & \swdA & \swdA & \swdA & No tiene \\
    Path cost hasta\\ el Root bridge & 19 & 19 & 25 & --- \\
    Root port & Fa0/3 & Fa0/3 & Gi0/1 & No tiene \\
    \hline
  \end{tabular}
  \captionof{table}{Informacion de cada switch acerca del Spanning Tree.}
  \label{tab:spanning-tree-before}
\end{center}

\begin{center}
  \begin{tabular}{ c c c c c }
    \hline
    --- & Switch-A & Switch-B & Switch-C & Switch-D \\
    \hline
    Switch-A & Alt & Alt & Alt & Designado \\
    Switch B & Alt & Alt & Alt & Designado \\
    Switch C & Alt & Alt & Alt & Designado \\
    Switch D & Root & Root & Root & --- \\
    PC (todas) & Designado & Designado & Designado & Designado \\
    \hline
  \end{tabular}
  \captionof{table}{Rol de cada switch. "Alt" es corto para "Alternativo".}
  \label{tab:spanning-tree-roles-before}
\end{center}

Una vez acabado de recopilar la informacion se paso a configurar el logging del switch
B con el proposito de analisar la \textit{conversacion} entre los switches
posteriormente. Esto se hizo usando los comandos mostrados en la terminal a
continuacion llamada "Switch B: Configurando logging", una vez se corrieron, tambien se
verifico la configuracion correcta usando \texttt{show debugging}.

Posteriormente se uso uno de los computadores conectados con el comando \texttt{ping -t
192.168.1.3}, se preparo un cronometro. Al desconectar el root bridge, se corrio el
tiempo desde el primer paquete fallido por el \texttt{ping} y se paro cuando el
\texttt{ping} volvio a recibir paquetes. El resultado fueron 24 segundos \textit{fuera
de linea}.

Finalmente se uso denuevo el mismo comando \texttt{show spanning-tree} para ver como
cambiaron los roles de red. Los resultados, denuevo se dejan en la tabla
\ref{tab:spanning-tree-after} y \ref{tab:spanning-tree-roles-after}. En estas se
evidencia como el Switch C paso de ser el alternativo al root a falta del switch D.


\begin{darkcode}[Switch B: Configurando logging]
Switch-B#debug spanning-tree events
'Spanning Tree event debugging is on'
Switch-B#debug spanning-tree events root
'Spanning Tree root changes debugging is on'
Switch-B#debug spanning-tree root pvst+
'Spanning Tree PVST+ debugging is on'
Switch-B#debug spanning-tree pvst+ switch state
'Spanning Tree Port state changes debugging is on'
Switch-B#debug pm vp
'Virtual port events debugging is on'
Switch-B#debug spanning-tree bpdu receive
'Spanning Tree BPDU Received debugging is on'
Switch-B#debug spanning-tree bpdu receive transmit
'Spanning Tree BPDU Transmitted debugging is on'
Switch-B#debug spanning-tree bpdu
'Spanning Tree Switch Shim process receive bpdu debugging is on'
Switch-B#debug spanning-tree switch rx process decode
'Spanning Tree Switch Shim decode received packets debugging is on'
Switch-B#show debugging
'
Spanning Tree:
  Spanning Tree BPDU Transmitted debugging is on
  Spanning Tree BPDU Received debugging is on
  Spanning Tree event debugging is on
  Spanning Tree root changes debugging is on
  Spanning Tree PVST+ debugging is on
Spanning Tree Switch Shim:
  Spanning Tree Port state changes debugging is on
  Spanning Tree Switch Shim decode received packets debugging is on
  Spanning Tree Switch Shim process receive bpdu debugging is on
Port Manager:
  Virtual port events debugging is on
\end{darkcode}

\begin{center}
  \begin{tabular}{ l c c c c }
    \hline
    --- & Switch-A & Switch-B & Switch-C & Switch-D \\
    \hline
    MAC & \swaA & \swbA & \swcA & --- \\
    BID & \swaA & \swbA & \swcA & --- \\
    Root bridge & \swcA & \swcA & \swcA & --- \\
    Path cost hasta\\ el Root bridge & 19 & 19 & --- & --- \\
    Root port & Fa0/2 & Fa0/2 & No tiene & --- \\
    \hline
  \end{tabular}
  \captionof{table}{Informacion de cada switch acerca del Spanning Tree.}
  \label{tab:spanning-tree-after}
\end{center}

\begin{center}
  \begin{tabular}{ c c c c c }
    \hline
    --- & Switch-A & Switch-B & Switch-C \\
    \hline
    Switch-A & Alt & Alt & Designado \\
    Switch B & Alt & Alt & Designado \\
    Switch C & Root & Root & --- \\
    PC (todas) & Designado & Designado & Designado \\
    \hline
  \end{tabular}
  \captionof{table}{Rol de cada switch. "Alt" es corto para "Alternativo".}
  \label{tab:spanning-tree-roles-after}
\end{center}

\subsubsection*{Analisis de los archivos de logging}
El archivo es demasiado largo para dejarsele en su totalidad, por lo que se resaltaran
lineas claves, esto incluye la primeria linea que hace referencia al root switch
estando caido:
\begin{center}
  [23]\texttt{*Mar  1 01:53:44.870: *** vp\_linkchange: single: \textcolor{red}{down}: 1/3(1) }
\end{center}
La cual nos indica que el primer tiempo en el que la red vio el fallo fue en
\texttt{53:44} (en formato \texttt{MM:SS}). Al principio, antes de apagar el switch D,
se veia la siguiente linea
\begin{center}
  [9] \texttt{*Mar  1 01:53:42.999:     cfg p:0000 v:00 t:00 f:00 r:8001 \textcolor{red}{001c.f6b4.2480} ...}
\end{center}
La direccion MAC resaltada se reconoce como la del switch D, el switch root. Despues de
que este mismo se apague sale esta linea inmediatamente
\begin{center}
  [19]\texttt{*Mar  1 01:53:44.870: STP: VLAN0001 new root port Fa0/2, cost 23}
\end{center}
Esto significa que el switch C (el que esta en tal puerto) es el nuevo root casi
inmediatamente despues de apagar el switch 2 segundos antes. Pero aun no esta
funcionando el ping por que el protocolo STP aun requeriere que esperar 15 segundos en
listening y 15 segundos en learning. Una vez pasado al rededor de ello (~24 segundos
despues de desconectar el switch D), se ve el ping funcionando denuevo y el switch B ya
no esta llamando ningun cambio de topologia como 

\noindent[277]\texttt{*Mar  1 01:54:02.998: STP SW: Gi0/1 new blocking req for 1 vlans}

\noindent[278]\texttt{*Mar  1 01:54:02.998: STP[1]: Generating TC trap for port GigabitEt...}

\noindent[279]\texttt{*Mar  1 01:54:02.998: STP: VLAN0001 Gi0/1 -> blocking}
